\begin{tikzpicture}[
  font=\small,
  >={Stealth[length=2.2mm]},
  data/.style={draw=teal!65!black,rounded corners=2pt,fill=teal!7,
    minimum height=15mm,text width=30mm,align=center,inner sep=2.5pt},
  process/.style={draw=blue!65!black,rounded corners=2pt,fill=blue!7,
    minimum height=15mm,text width=30mm,align=center,inner sep=2.5pt},
  robot/.style={draw=orange!75!black,rounded corners=2pt,fill=orange!10,
    minimum height=15mm,text width=30mm,align=center,inner sep=2.5pt},
  decide/.style={draw=red!55!black,rounded corners=2pt,fill=red!4,
    minimum height=15mm,text width=30mm,align=center,inner sep=2.5pt},
  recover/.style={draw=violet!65!black,rounded corners=2pt,fill=violet!7,
    minimum height=15mm,text width=30mm,align=center,inner sep=2.5pt},
  calib/.style={draw=gray!65,rounded corners=5pt,fill=gray!8,
    text width=30mm,align=center,inner sep=2.5pt,font=\footnotesize},
  arrow/.style={->,thick,blue!70!black},
  dataarrow/.style={->,thick,teal!70!black},
  alarm/.style={->,thick,red!65!black},
  feedback/.style={->,thick,orange!80!black,rounded corners=3pt},
  influence/.style={->,densely dashed,gray!70,rounded corners=2pt},
  databus/.style={thick,teal!70!black}
]
  \node[data] (prev) {\textbf{前一RTC动作块}\\时间对齐的\\未执行尾部};
  \node[process,right=6mm of prev] (rtc) {\textbf{RTC连续衔接}\\冻结\(d_i\)步前缀\\引导重叠区并补全};
  \node[robot,right=6mm of rtc] (exec) {\textbf{执行与观测}\\执行RTC动作块\\获取新观测};
  \node[data,below=8mm of prev] (raw) {\textbf{当前原始提议}\\新观测+\(\pi_{0.5}\)\\固定参考噪声};
  \node[process,below=8mm of rtc] (inc) {\textbf{跨块一致性残差}\\归一化\(r_i\)\\滑动均值\(\widetilde r_i\)};
  \node[decide,below=8mm of exec] (dec) {\textbf{报警判定}\\\(\widetilde r_i>\eta\)\\连续\(n_{\mathrm{hold}}\)次？};
  \node[recover,below=8mm of dec] (rec) {\textbf{模板恢复}\\停止当前块\\退至\(z_{\mathrm{safe}}\)→回入口位姿};
  \node[calib,below=8mm of raw] (thr) {\textbf{阈值标定（离线）}\\成功回合的全程最大值\\取95分位};
  \node[calib,below=8mm of inc] (shadow) {\textbf{离线检测评价}\\回放记录的C2回合\\只记录报警，不改变动作};
  \node[recover,below=8mm of rec] (retry) {\textbf{重试当前子任务一次}\\成功则继续后序任务\\整链不再允许第二次恢复};
  \node[decide,below=8mm of retry] (stop) {\textbf{安全停机}\\条件不满足或再次报警\\交人工处理};

  % 两份时间对齐的动作提议先在节点间的空白带汇入数据总线，
  % 再分别送往RTC衔接与一致性计算，避免交叉连线穿过文本块。
  \coordinate (bus-top) at ($(prev.east)!0.5!(rtc.west)$);
  \coordinate (bus-bottom) at ($(raw.east)!0.5!(inc.west)$);
  \draw[databus] (prev.east) -- (bus-top) -- (bus-bottom) -- (raw.east);
  \draw[dataarrow] (bus-top) -- (rtc.west);
  \draw[dataarrow] (bus-bottom) -- (inc.west);
  \draw[arrow] (rtc) -- (exec);
  \draw[arrow] (inc) -- (dec);
  % 离线信息只从节点间通道接入，不与在线恢复主链共线。
  \draw[influence] (thr.north) -- ++(0,4mm) -|
    node[pos=0.34,above,font=\footnotesize,fill=white,inner sep=1pt]{既定阈值}
    ($(dec.south west)+(5mm,0)$);
  \draw[influence] (inc.south) -- (shadow.north);
  \draw[arrow] (dec.north) -- node[right,font=\footnotesize]{否：继续} (exec.south);
  \draw[alarm] (dec) -- node[right,font=\footnotesize,fill=white,inner sep=1pt]{C3在线：是} (rec);
  \draw[arrow] (rec) -- (retry);
  \draw[alarm] (retry) -- node[right,font=\footnotesize]{不安全/再次报警} (stop);
  \draw[feedback] (retry.east) -- ++(6mm,0) |-
    node[pos=0.24,right,font=\footnotesize,align=left,fill=white,inner sep=1pt]{恢复后\\重新推理}
    (exec.east);
  \draw[feedback] (exec.north) -- ++(0,6mm) -|
    node[pos=0.24,above,font=\footnotesize,fill=white,inner sep=1pt]{新观测→下一次\(\pi_{0.5}\)推理}
    ([xshift=-7mm]raw.west) -- (raw.west);
  \node[draw=gray!45,densely dashed,rounded corners=3pt,fit=(thr)(shadow),
    inner sep=2.5mm,label={[font=\footnotesize,text=gray!70]below:离线标定与检测评价}] {};
\end{tikzpicture}
